
\documentclass{article}
\usepackage{colortbl}
\usepackage{makecell}
\usepackage{multirow}
\usepackage{supertabular}

\begin{document}

\newcounter{utterance}

\centering \large Interaction Transcript for game `dond', experiment `coop\_en', episode 10 with reflexion{-}Qwen3.5{-}27B.
\vspace{24pt}

{ \footnotesize  \setcounter{utterance}{1}
\setlength{\tabcolsep}{0pt}
\begin{supertabular}{c@{$\;$}|p{.15\linewidth}@{}p{.15\linewidth}p{.15\linewidth}p{.15\linewidth}p{.15\linewidth}p{.15\linewidth}}
    \# & $\;$A & \multicolumn{4}{c}{Game Master} & $\;\:$B\\
    \hline

    \theutterance \stepcounter{utterance}  
    & & \multicolumn{4}{p{0.6\linewidth}}{
        \cellcolor[rgb]{0.9,0.9,0.9}{
            \makecell[{{p{\linewidth}}}]{
                \texttt{\tiny{[P1$\langle$GM]}}
                \texttt{You are playing a cooperative negotiation game in which you have to agree on how to divide a set of items among you and another player.} \\
\\ 
\texttt{Rules:} \\
\texttt{(a) You and the other player are given a set of items. Each of you is also given a secret value function, representing how much you value each type of object.} \\
\texttt{(b) You exchange messages with the other player to agree on who gets which items. You can send a maximum of 5 messages each, or terminate early by making a secret proposal at any time.} \\
\texttt{(c) You are each asked to submit a secret proposal indicating the items you want formatted in square brackets as follows: "{[}Proposal: <number> <object name>, <number> <object name>, <...>{]}"} \\
\texttt{(d) If your proposals are complementary, i.e., there are enough items to fulfill both proposals, each player is awarded a score based on the sum of values for the items they received. Otherwise, both of you get zero points.} \\
\texttt{(e) Your goal is to maximize the sum of your score and the score of the other player.} \\
\\ 
\texttt{Let us start.} \\
\\ 
\texttt{The set of available items is:} \\
\\ 
\texttt{1 glass, 2 wires, 2 chalks} \\
\\ 
\texttt{Your secret value function is:} \\
\\ 
\texttt{glass: 8, wire: 1, chalk: 0} \\
\\ 
\texttt{IMPORTANT: Your messages, unless it is the secret proposal, are directly transmitted to the other player, so do not include any response to the rules or text announcing your message. To make a secret proposal, use the indicated format. Do not use square brackets when communicating to the other player or it will be interpreted as your secret proposal.} \\
            }
        }
    }
    & \\ \\

    \theutterance \stepcounter{utterance}  
    & \multicolumn{4}{p{0.6\linewidth}}{
        \cellcolor[rgb]{0.9,0.9,0.9}{
            \makecell[{{p{\linewidth}}}]{
                \texttt{\tiny{[P1$\rangle$GM]}}
                \texttt{I value the glass highly (8 points) and don't care about the wires or chalks. Since we want to maximize our combined score, it makes sense for me to take the glass. You can have the 2 wires and 2 chalks. Does that work for you?} \\
            }
        }
    }
    & & \\ \\

    \theutterance \stepcounter{utterance}  
    & & & \multicolumn{2}{p{0.3\linewidth}}{
        \cellcolor[rgb]{0.9,0.9,0.9}{
            \makecell[{{p{\linewidth}}}]{
                \texttt{\tiny{[GM$|$GM]}}
                \texttt{continue} \\
            }
        }
    }
    & & \\ \\

    \theutterance \stepcounter{utterance}  
    & & & \multicolumn{2}{p{0.3\linewidth}}{
        \cellcolor[rgb]{0.9,0.9,0.9}{
            \makecell[{{p{\linewidth}}}]{
                \texttt{\tiny{[GM$|$GM]}}
                \texttt{I value the glass highly (8 points) and don't care about the wires or chalks. Since we want to maximize our combined score, it makes sense for me to take the glass. You can have the 2 wires and 2 chalks. Does that work for you?} \\
            }
        }
    }
    & & \\ \\

    \theutterance \stepcounter{utterance}  
    & & \multicolumn{4}{p{0.6\linewidth}}{
        \cellcolor[rgb]{0.9,0.9,0.9}{
            \makecell[{{p{\linewidth}}}]{
                \texttt{\tiny{[GM$\rangle$P2]}}
                \texttt{You are playing a cooperative negotiation game in which you have to agree on how to divide a set of items among you and another player.} \\
\\ 
\texttt{Rules:} \\
\texttt{(a) You and the other player are given a set of items. Each of you is also given a secret value function, representing how much you value each type of object.} \\
\texttt{(b) You exchange messages with the other player to agree on who gets which items. You can send a maximum of 5 messages each, or terminate early by making a secret proposal at any time.} \\
\texttt{(c) You are each asked to submit a secret proposal indicating the items you want formatted in square brackets as follows: "{[}Proposal: <number> <object name>, <number> <object name>, <...>{]}"} \\
\texttt{(d) If your proposals are complementary, i.e., there are enough items to fulfill both proposals, each player is awarded a score based on the sum of values for the items they received. Otherwise, both of you get zero points.} \\
\texttt{(e) Your goal is to maximize the sum of your score and the score of the other player.} \\
\\ 
\texttt{Let us start.} \\
\\ 
\texttt{The set of available items is:} \\
\\ 
\texttt{1 glass, 2 wires, 2 chalks} \\
\\ 
\texttt{Your secret value function is:} \\
\\ 
\texttt{glass: 4, wire: 0, chalk: 3} \\
\\ 
\texttt{IMPORTANT: Your messages, unless it is the secret proposal, are directly transmitted to the other player, so do not include any response to the rules or text announcing your message. To make a secret proposal, use the indicated format. Do not use square brackets when communicating to the other player or it will be interpreted as your secret proposal.} \\
\\ 
\\ 
\texttt{I value the glass highly (8 points) and don't care about the wires or chalks. Since we want to maximize our combined score, it makes sense for me to take the glass. You can have the 2 wires and 2 chalks. Does that work for you?} \\
            }
        }
    }
    & \\ \\

    \theutterance \stepcounter{utterance}  
    & & & \multicolumn{4}{p{0.6\linewidth}}{
        \cellcolor[rgb]{0.9,0.9,0.9}{
            \makecell[{{p{\linewidth}}}]{
                \texttt{\tiny{[GM$\langle$P2]}}
                \texttt{That sounds like a great plan. Since I value the glass at 4 points and the chalks at 3 points each (totaling 6 for the chalks), while I get 0 for the wires, I would prefer to keep the chalks and give you the glass.} \\
\\ 
\texttt{However, your proposal gives me the wires (which I value at 0) and the chalks. If I take the 2 chalks, my score is 6. If you take the glass, your score is 8. The total combined score would be 14.} \\
\\ 
\texttt{Wait, let me re{-}evaluate based on my values:} \\
\texttt{{-} Glass: 4} \\
\texttt{{-} Wire: 0} \\
\texttt{{-} Chalk: 3} \\
\\ 
\texttt{If I take the 2 chalks, I get 6 points.} \\
\texttt{If you take the glass, you get 8 points.} \\
\texttt{Total: 14.} \\
\\ 
\texttt{If I took the glass instead, I would get 4 points, and you would get 0 (since you don't care about wires or chalks? Wait, you said you value glass highly at 8, but didn't mention wires or chalks. Assuming you value them at 0 or low).} \\
\\ 
\texttt{Actually, looking at your message again: "I value the glass highly (8 points) and don't care about the wires or chalks."} \\
\texttt{My values: Glass 4, Wire 0, Chalk 3.} \\
\\ 
\texttt{Option A (Your proposal):} \\
\texttt{{-} You: Glass (8)} \\
\texttt{{-} Me: 2 Wires (0) + 2 Chalks (6) = 6} \\
\texttt{{-} Total: 14} \\
\\ 
\texttt{Option B (Alternative):} \\
\texttt{{-} You: Glass (8)} \\
\texttt{{-} Me: 2 Chalks (6) + 2 Wires (0) = 6} \\
\texttt{{-} Total: 14} \\
\\ 
\texttt{It seems the distribution of wires doesn't matter for the score since we both value them at 0. The key is ensuring I get the chalks and you get the glass.} \\
\\ 
\texttt{Your proposal: "You can have the 2 wires and 2 chalks."} \\
\texttt{This works perfectly for me. I will take the 2 chalks and the 2 wires (even though wires are worthless to me, it doesn't hurt). You take the glass.} \\
\\ 
\texttt{Let's finalize this.} \\
\\ 
\texttt{{[}Proposal: 1 glass{]}} \\
            }
        }
    }
     \\ \\

    \theutterance \stepcounter{utterance}  
    & & & \multicolumn{2}{p{0.3\linewidth}}{
        \cellcolor[rgb]{0.9,0.9,0.9}{
            \makecell[{{p{\linewidth}}}]{
                \texttt{\tiny{[GM$|$GM]}}
                \texttt{proposal} \\
            }
        }
    }
    & & \\ \\

    \theutterance \stepcounter{utterance}  
    & & & \multicolumn{2}{p{0.3\linewidth}}{
        \cellcolor[rgb]{0.9,0.9,0.9}{
            \makecell[{{p{\linewidth}}}]{
                \texttt{\tiny{[GM$|$GM]}}
                [1, 0, 0]
            }
        }
    }
    & & \\ \\

    \theutterance \stepcounter{utterance}  
    & & \multicolumn{4}{p{0.6\linewidth}}{
        \cellcolor[rgb]{0.9,0.9,0.9}{
            \makecell[{{p{\linewidth}}}]{
                \texttt{\tiny{[P1$\langle$GM]}}
                \texttt{The other player has submitted a secret proposal. It is now your turn to submit a proposal of your own.} \\
            }
        }
    }
    & \\ \\

    \theutterance \stepcounter{utterance}  
    & \multicolumn{4}{p{0.6\linewidth}}{
        \cellcolor[rgb]{0.9,0.9,0.9}{
            \makecell[{{p{\linewidth}}}]{
                \texttt{\tiny{[P1$\rangle$GM]}}
                \texttt{{[}Proposal: 1 glass{]}} \\
            }
        }
    }
    & & \\ \\

    \theutterance \stepcounter{utterance}  
    & & & \multicolumn{2}{p{0.3\linewidth}}{
        \cellcolor[rgb]{0.9,0.9,0.9}{
            \makecell[{{p{\linewidth}}}]{
                \texttt{\tiny{[GM$|$GM]}}
                \texttt{proposal} \\
            }
        }
    }
    & & \\ \\

    \theutterance \stepcounter{utterance}  
    & & & \multicolumn{2}{p{0.3\linewidth}}{
        \cellcolor[rgb]{0.9,0.9,0.9}{
            \makecell[{{p{\linewidth}}}]{
                \texttt{\tiny{[GM$|$GM]}}
                [1, 0, 0]
            }
        }
    }
    & & \\ \\

    \theutterance \stepcounter{utterance}  
    & & & \multicolumn{2}{p{0.3\linewidth}}{
        \cellcolor[rgb]{0.9,0.9,0.9}{
            \makecell[{{p{\linewidth}}}]{
                \texttt{\tiny{[GM$|$GM]}}
                [[1, 0, 0], [1, 0, 0]]
            }
        }
    }
    & & \\ \\

    \theutterance \stepcounter{utterance}  
    & \multicolumn{4}{p{0.6\linewidth}}{
        \cellcolor[rgb]{0.95,0.95,0.95}{
            \makecell[{{p{\linewidth}}}]{
                \texttt{\tiny{[P2$\rightarrow$P1]}}
                \texttt{Post{-}episode Reflexion (not sent to GM):} \\
\texttt{New plan: I failed because I assumed the other player would accept my initial verbal offer and submit a complementary proposal, but they submitted their own secret proposal first, likely claiming items I didn't explicitly secure or creating a conflict. My strategy of making a unilateral verbal offer without confirming the other player's specific valuation or securing a mutual agreement before the deadline was flawed. Next time, I will explicitly ask the other player for their secret value function or their specific item preferences before making a proposal. I will then negotiate a specific split (e.g., "I take the glass, you take the wires and chalks") and wait for their explicit confirmation ("Yes, that works") before submitting my secret proposal. If they submit a proposal first, I must immediately analyze it against the total inventory to ensure my counter{-}proposal is strictly complementary, rather than assuming my previous verbal offer was binding.} \\
            }
        }
    }
    & & \\ \\

    \theutterance \stepcounter{utterance}  
    & & & \multicolumn{4}{p{0.6\linewidth}}{
        \cellcolor[rgb]{0.95,0.95,0.95}{
            \makecell[{{p{\linewidth}}}]{
                \texttt{\tiny{[P1$\rightarrow$P2]}}
                \texttt{Post{-}episode Reflexion (not sent to GM):} \\
\texttt{New plan: I failed because I submitted a proposal claiming only the glass (`{[}Proposal: 1 glass{]}`) while the other player's proposal (implied by the negotiation) likely claimed the chalks and wires. Since the rules state that proposals must be *complementary* (covering all items without overlap) to score, leaving the chalks and wires unclaimed in my proposal caused a mismatch, resulting in zero points for both. I mistakenly assumed that agreeing verbally was sufficient or that the other player would fill the gaps, but the game requires my specific proposal to explicitly list *all* items I am claiming. Next time, after agreeing on the split (Me: 2 chalks, 2 wires; You: 1 glass), I will submit a proposal that explicitly lists every item I am taking: `{[}Proposal: 2 chalks, 2 wires{]}`. This ensures the union of our proposals covers the entire set \{1 glass, 2 wires, 2 chalks\} exactly once.} \\
            }
        }
    }
     \\ \\

\end{supertabular}
}

\end{document}
